% Vitesse d'un parachutiste en chute libre avec frottement quadratique (résultat du programme, chapitre 14)
\begin{tikzpicture}
\begin{axis}[otMesures,otAxesBout,otLecture,width=98mm,height=44mm,xmin=0,xmax=31,ymin=0,ymax=62,
  xtick={0,5,...,30},ytick={0,10,...,60},
  xlabel={$t$ (s)},ylabel={$v$ (m$\cdot$s$^{-1}$)}]
\addplot[otMuted,line width=0.6pt,dash pattern=on 2pt off 1.5pt,domain=0:5.5,samples=2] {9.81*x};
\addplot[otOrange,line width=0.6pt,dash pattern=on 2pt off 1.5pt,domain=0:30,samples=2] {53.9};
\addplot[otNavy,line width=1pt,domain=0:30,samples=150] {53.9*tanh(9.81*x/53.9)};
\node[pcLegende,text=otOrange!80!black,anchor=south east] at (axis cs:30,54.5) {vitesse limite : \qty{53,9}{m.s^{-1}}};
\node[pcLegende,text=otMuted,anchor=west,rotate=0] at (axis cs:5.6,50) {chute sans frottement};
\end{axis}
\end{tikzpicture}
\par\smallskip\legende{Au début, la vitesse croît comme en chute libre ; le frottement augmente avec la vitesse et finit par compenser le poids.}
